%% toc: P2 20 20.5,20.6
\section{Proofs on the set-theoretic prerequisites of the L\"owenheim--Skolem theorems}
\label{sec:annex-p02-w02b-p02-029-ls-prerequisites}

Notation is that of \cref{sec:p02-w02b-p02-029-ls-prerequisites}. References
to items of the ordinal lemma, the finite-set lemma and the infinite-ordinal
lemma are to \cref{thm:p01-w02-p01-011-ordinals},
\cref{thm:p01-w02-p01-011-finite-sets} and
\cref{thm:p01-w02-p01-011-infinite-ordinals}; ``cardinal assignment'' is
\cref{thm:p01-w02-p01-008-cardinal-assignment}; ``CSB'' is the
Cantor--Schr\"oder--Bernstein theorem of \cref{sec:set-prerequisites-05}.
Induction on \(\N\) is \cref{thm:p01-w02-p01-008-induction}.

\begin{annexitem}{annex:p02-w02b-p02-029-bounds}{thm:p02-w02b-p02-029-bounds}{Closure of the bound \(\abs X\leq\mu\)}
\begin{proof}
Fix injections \(e:X\to\mu\) and \(e':Y\to\mu\); they exist by hypothesis.
By item~1 of the infinite-ordinal lemma, \(\N\subseteq\mu\); in particular
\(0,1\in\mu\). By \cref{thm:p01-w02-p01-011-square}, fix a bijection
\(\pi:\mu\times\mu\to\mu\).

\emph{Item~1.} The restriction of \(e\) to \(Z\) is an injection
\(Z\to\mu\).

\emph{Item~2.} If \(Z\) is countable, there is an injection \(Z\to\N\), and
its composite with the inclusion \(\N\subseteq\mu\) is an injection
\(Z\to\mu\). A finite set is countable by item~6 of
\cref{thm:p01-w02-p01-008-pigeonhole}.

\emph{Item~3.} The map \((x,y)\mapsto\pi(e(x),e'(y))\) is an injection
\(X\times Y\to\mu\): equal values give equal pairs \((e(x),e'(y))\) because
\(\pi\) is injective, hence \(x\) and \(y\) are determined because \(e\) and
\(e'\) are injective. For the union, define \(j\) on \(X\cup Y\) by
\(j(z)=\pi(0,e(z))\) if \(z\in X\) and \(j(z)=\pi(1,e'(z))\) if
\(z\in Y\setminus X\); this is a set by separation from
\((X\cup Y)\times\mu\). If \(j(z)=j(z')\), injectivity of \(\pi\) gives the
same first coordinate. If it is \(0\), both lie in \(X\) and
\(e(z)=e(z')\), so \(z=z'\); if it is \(1\), both lie in \(Y\setminus X\) and
\(e'(z)=e'(z')\), so \(z=z'\). So \(j\) is an injection
\(X\cup Y\to\mu\).

\emph{Item~4.} For \(w\in g[X]\), the set
\(S_w=\{e(x):x\in X,\ g(x)=w\}\) is a nonempty set of ordinals, being a
subset of \(\mu\) whose members are ordinals by item~1 of the ordinal lemma.
By item~4 of the ordinal lemma it has a unique least member; let \(h(w)\) be
that member. By separation from \(g[X]\times\mu\), \(h\) is a function
\(g[X]\to\mu\). If \(h(w)=h(w')\), then \(h(w)=e(x)\) with \(g(x)=w\) and
\(h(w')=e(x')\) with \(g(x')=w'\); injectivity of \(e\) gives \(x=x'\), and
then \(w=g(x)=g(x')=w'\). So \(h\) is an injection.

\emph{Item~5.} This is item~1 of
\cref{thm:p01-w02-p01-011-syntax-counting} with \(\kappa=\mu\) and
\(A=X\).

In items~1--4 the functions \(e\), \(e'\) and \(\pi\) are single
instantiations of existence hypotheses inside the proof of one existence
statement, and \(h\) is defined by separation from them; no choice principle
is used. Item~5 inherits the use of the Axiom of Choice in clause~2 of
\cref{thm:p01-w02-p01-011-finite-coding}.
\end{proof}
\end{annexitem}

\begin{annexitem}{annex:p02-w02b-p02-029-mu}{thm:p02-w02b-p02-029-mu}{The cardinal of the downward bound}
\begin{proof}
By items~1 and~2 of cardinal assignment, \(A\) and \(\mathcal F\cup\mathcal R\)
are each bijective with exactly one cardinal, \(\abs A\) and
\(\abs\Lang\) respectively; by item~5, \(\aleph_0=\N\) is the cardinal of
\(\N\). By item~4, the list \(\abs A,\abs\Lang,\aleph_0\) has a greatest
member \(\mu\); it is one of the three, so it is a cardinal, and each of the
three is \(\leq\mu\) in the ordinal order.

For ordinals, \(\nu\leq\mu\) iff \(\nu\subseteq\mu\) (item~3 of the ordinal
lemma). So \(\abs A\subseteq\mu\), and composing a bijection
\(A\to\abs A\) with the inclusion gives an injection \(A\to\mu\); likewise
for \(\mathcal F\cup\mathcal R\), and \(\N\subseteq\mu\) gives the inclusion
\(\N\to\mu\).

Suppose \(\mu\) were bijective with some \(n\in\N\), by \(b:\mu\to n\). Since
\(\N\subseteq\mu\), the restriction of \(b\) to \(\N\) is an injection
\(\N\to n\), which item~4 of \cref{thm:p01-w02-p01-008-pigeonhole} excludes.
So \(\mu\) is an infinite cardinal.
\end{proof}
\end{annexitem}

\begin{annexitem}{annex:p02-w02b-p02-029-finite-injections}{thm:p02-w02b-p02-029-finite-injections}{Finite pieces of an infinite set}
\begin{proof}
We prove by induction on \(n\in\N\) that an injection \(n\to D\) exists. For
\(n=0=\varnothing\), the empty function is one. Let \(e:n\to D\) be an
injection. It is not surjective: otherwise it would be a bijection
\(n\to D\), and \(D\) is bijective with no natural number. So there is
\(d\in D\setminus e[n]\). Since \(n\notin n\) (item~6 of the ordinal lemma),
\(e\cup\{(n,d)\}\) is a function on \(n\cup\{n\}=n+1\). It is injective:
\(e\) is injective, and its value \(d\) at \(n\) is not a value of \(e\).
This proves the induction step.

If \(I\) is finite, there are \(m\in\N\) and a bijection \(b:m\to I\). With
an injection \(e:m\to D\), the composite \(e\circ b^{-1}\) is an injection
\(I\to D\).

Each step instantiates one existence statement (the injection \(e\), the
element \(d\)) in proving another existence statement, so the induction
proves, for each \(n\), the existence of an injection; it does not select a
sequence of injections, and no choice principle is used.
\end{proof}
\end{annexitem}

\begin{annexitem}{annex:p02-w02b-p02-029-skolem-functions}{thm:p02-w02b-p02-029-skolem-functions}{Witness functions}
\begin{proof}
Write \(D=\dom\mathcal M\), a nonempty set. By the well-ordering theorem
(\cref{thm:p01-w02-p01-008-well-ordering}), which uses the Axiom of Choice
once, fix a relation \(\prec\) that strictly well-orders \(D\). Nothing else
below uses a choice principle.

For a pattern \(e=(\varphi,(y,x_1,\ldots,x_k))\) and
\(\bar a\in D^k\), let
\(W_e(\bar a)=\{b\in D:\mathcal M\entails\varphi(b,\bar a)\}\), a set by
separation from \(D\). If \(W_e(\bar a)\neq\varnothing\), it has a
\(\prec\)-least member, and only one: two least members \(b,b'\) would
satisfy neither \(b\prec b'\) nor \(b'\prec b\), so \(b=b'\) because
\(\prec\) is a strict linear order. By separation from
\(\mathcal E(\Lang)\times D^{<\omega}\times D\), let \(\Phi\) be the set of
triples \((e,\bar a,b)\) such that \(\bar a\) has length equal to the arity
of \(e\), and either \(b\) is the \(\prec\)-least member of a nonempty
\(W_e(\bar a)\), or \(W_e(\bar a)=\varnothing\) and \(b=m_0\). For each
\(e\) and each \(\bar a\) of the right length there is exactly one such
\(b\). By replacement, \(e\mapsto F_e=\{(\bar a,b):(e,\bar a,b)\in\Phi\}\) is a
function on \(\mathcal E(\Lang)\), and \(F_e\) is a function
\(D^k\to D\) for \(e\) of arity \(k\).

If \(\mathcal M\entails\exists y\,\varphi(y,\bar a)\), the satisfaction
clause for \(\exists\) (\cref{def:p02-w02b-p02-029-patterns}) gives some
\(b\in D\) with \(\mathcal M\entails\varphi(b,\bar a)\); so
\(W_e(\bar a)\neq\varnothing\), and \(F_e(\bar a)\), its least member, lies
in \(W_e(\bar a)\): \(\mathcal M\entails\varphi(F_e(\bar a),\bar a)\). This
is item~1. Otherwise no such \(b\) exists, \(W_e(\bar a)=\varnothing\), and
\(F_e(\bar a)=m_0\). This is item~2.
\end{proof}
\end{annexitem}

\begin{annexitem}{annex:p02-w02b-p02-029-stages}{thm:p02-w02b-p02-029-stages}{Stage bounds and closure}
\begin{proof}
Write \(D=\dom\mathcal M\). By \cref{thm:p02-w02b-p02-029-mu}, \(\mu\) is an
infinite cardinal with \(\abs A\leq\mu\) and
\(\abs{\mathcal F\cup\mathcal R}\leq\mu\). Items below of the form
``item~\(i\)'' without further reference are items of
\cref{thm:p02-w02b-p02-029-bounds}.

\emph{Counting the indices.} \(\abs{\mathcal F}\leq\mu\) by item~1. Let
\(\Gamma=\mathcal F\cup\mathcal R\cup\mathsf{Var}\cup T_{\mathrm{log}}\) be the
alphabet of \(\Lang\). The set \(\mathsf{Var}\) is countable and
\(T_{\mathrm{log}}\) has eleven members, so both are \(\leq\mu\) by item~2,
and \(\abs\Gamma\leq\mu\) by item~3 applied twice. Every formula is a string
over \(\Gamma\) (\cref{sec:fol-syntax-semantics-02}), that is, a member of
\(\Gamma^{<\omega}\), the set of finite sequences of \cref{sec:strings-01};
so \(\abs{\mathsf{Fm}(\Lang)}\leq\mu\) by items~5 and~1. This is the bound of
item~2 of \cref{thm:p01-w02-p01-011-syntax-counting}. Also
\(\abs{\mathsf{Var}^{<\omega}}\leq\mu\) by item~5, and
\(\mathcal E(\Lang)\subseteq\mathsf{Fm}(\Lang)\times\mathsf{Var}^{<\omega}\),
so \(\abs{\mathcal E(\Lang)}\leq\mu\) by items~3 and~1.

\emph{Item~1.} Let \(X\subseteq D\) with \(\abs X\leq\mu\). Then
\(\abs{X^{<\omega}}\leq\mu\) by item~5. Let
\(P=\{(f,\bar a)\in\mathcal F\times X^{<\omega}:f\text{ has arity }k\geq1
\text{ and }\bar a\in X^k\}\). By items~3 and~1, \(\abs P\leq\mu\), and
\(\operatorname{Fn}(X)\) is the image of \(P\) under
\((f,\bar a)\mapsto f^{\mathcal M}(\bar a)\), so
\(\abs{\operatorname{Fn}(X)}\leq\mu\) by item~4. Let
\(Q=\{(e,\bar a)\in\mathcal E(\Lang)\times X^{<\omega}:\bar a\in X^k\text{
for the arity }k\text{ of }e\}\). Again \(\abs Q\leq\mu\), and
\(\operatorname{Wt}(X)\) is the image of \(Q\) under
\((e,\bar a)\mapsto F_e(\bar a)\), which is defined because
\(X^k\subseteq D^k\); so \(\abs{\operatorname{Wt}(X)}\leq\mu\). Item~3
applied twice gives \(\abs{\operatorname{St}(X)}\leq\mu\). Moreover
\(\operatorname{St}(X)\subseteq D\), since \(f^{\mathcal M}\) and \(F_e\)
take values in \(D\).

\emph{Item~2.} \(A_0\subseteq D\). Its three parts satisfy the bound:
\(\abs A\leq\mu\); \(\{m_0\}\) is finite, so item~2 applies; and the set of
values \(c^{\mathcal M}\) is the image of the set of constant symbols, a
subset of \(\mathcal F\), so items~1 and~4 apply. By item~3,
\(\abs{A_0}\leq\mu\). By induction on \(n\), using item~1 of the present
theorem for the step, \(A_n\subseteq D\) and \(\abs{A_n}\leq\mu\) for every
\(n\in\N\). The hypothesis of clause~4 of
\cref{thm:p01-w02-p01-011-finite-coding} with \(\kappa=\mu\) is therefore
satisfied, and that clause gives \(\abs N\leq\mu\). This is where the
Axiom of Choice is needed: the induction proves for each \(n\) that an
injection \(A_n\to\mu\) exists, and the clause selects them all at once.

\emph{Item~3.} \(A_n\subseteq\operatorname{St}(A_n)=A_{n+1}\) for every
\(n\). We prove by induction on \(n\) that \(A_m\subseteq A_n\) for every
\(m\leq n\). For \(n=0\), \(m\leq0\) means \(m=0\). If the claim holds for
\(n\) and \(m\leq n+1\), then \(m\leq n\) or \(m=n+1\) by item~5 of the
ordinal lemma; in the first case \(A_m\subseteq A_n\subseteq A_{n+1}\). Next,
\(A\subseteq A_0\subseteq N\), and \(m_0\in A_0\subseteq N\), so
\(N\neq\varnothing\).

\emph{Item~4.} By induction on \(k\). The only member of \(N^0\) is the empty
sequence, which lies in \(A_0^0\). Let \(\bar a\in N^{k+1}\). By the
induction hypothesis, \(\bar a\mathbin{\upharpoonright}k\in A_{n'}^k\) for
some \(n'\), and \(\bar a(k)\in N\) lies in some \(A_{n''}\). By item~3 of
the ordinal lemma, \(n'\) and \(n''\) are comparable; let \(n\) be the larger.
By item~3 of the present theorem, \(A_{n'}\cup A_{n''}\subseteq A_n\), so
every coordinate of \(\bar a\) lies in \(A_n\), and \(\bar a\in A_n^{k+1}\).

\emph{Item~5.} Every \(c^{\mathcal M}\) lies in \(A_0\subseteq N\). Let \(f\)
have arity \(k\geq1\) and \(\bar a\in N^k\). By item~4, \(\bar a\in A_n^k\)
for some \(n\), so
\(f^{\mathcal M}(\bar a)\in\operatorname{Fn}(A_n)\subseteq A_{n+1}\subseteq N\).

\emph{Item~6.} Let \(\bar a\in N^k\) and
\(\mathcal M\entails\exists y\,\varphi(y,\bar a)\). By item~4,
\(\bar a\in A_n^k\) for some \(n\). Then
\(b=F_e(\bar a)\in\operatorname{Wt}(A_n)\subseteq A_{n+1}\subseteq N\), and
item~1 of \cref{thm:p02-w02b-p02-029-skolem-functions} gives
\(\mathcal M\entails\varphi(b,\bar a)\).
\end{proof}
\end{annexitem}

\begin{annexitem}{annex:p02-w02b-p02-029-fresh}{thm:p02-w02b-p02-029-fresh}{Fresh constant symbols}
\begin{proof}
By the union axiom applied twice, \(U=\bigcup\bigcup\Gamma\) is a set; by
separation, so is \(r=\{u\in U:u\notin u\}\). If \(r\in U\), then
\(r\in r\) iff \(r\notin r\), which is impossible; so \(r\notin U\). For
\(\alpha<\kappa\) let \(c_\alpha\) be the ordered pair
\((r,\alpha)=\{\{r\},\{r,\alpha\}\}\) of \cref{sec:set-prerequisites-02};
by replacement, \(\alpha\mapsto c_\alpha\) is a function on \(\kappa\). If
\(c_\alpha=c_\beta\), the ordered-pair identity of
\cref{sec:set-prerequisites-02} gives \(\alpha=\beta\). If
\(c_\alpha\in\Gamma\), then \(\{r\}\in c_\alpha\) gives
\(\{r\}\in\bigcup\Gamma\), and \(r\in\{r\}\) gives
\(r\in\bigcup\bigcup\Gamma=U\), which is false. So no \(c_\alpha\) is a
member of \(\Gamma\). The construction is explicit and uses no choice.
\end{proof}
\end{annexitem}

\begin{annexitem}{annex:p02-w02b-p02-029-upward-count}{thm:p02-w02b-p02-029-upward-count}{Cardinal steps of upward L\"owenheim--Skolem}
\begin{proof}
Let \(C=\{c_\alpha:\alpha<\kappa\}\); \(\alpha\mapsto c_\alpha\) is a
bijection \(\kappa\to C\). Since \(\aleph_0\leq\kappa\), the argument of the
last paragraph of the proof of \cref{thm:p02-w02b-p02-029-mu} shows that
\(\kappa\) is an infinite cardinal. Since \(\abs\Lang\leq\kappa\), item~3
of cardinal assignment gives an injection
\(\mathcal F\cup\mathcal R\to\kappa\).

\emph{Item~1.} The symbol set of \(\Lang^*\) is
\((\mathcal F\cup\mathcal R)\cup C\). Both parts are \(\leq\kappa\), so the
union is \(\leq\kappa\) by item~3 of \cref{thm:p02-w02b-p02-029-bounds}.
Conversely the inclusion of \(C\) composed with the bijection
\(\kappa\to C\) is an injection \(\kappa\to(\mathcal F\cup\mathcal R)\cup C\).
CSB gives a bijection, and item~2 of cardinal assignment gives
\(\abs{\Lang^*}=\kappa\).

\emph{Item~2.} First, if \(m\in\N\) and \(g\) is a function on \(m\), then
\(g[m]\) is finite: the proof of item~4 of
\cref{thm:p02-w02b-p02-029-bounds} with \(e\) the identity of \(m\) gives an
injection \(h:g[m]\to m\), so \(g[m]\) is bijective with \(h[g[m]]\), a
subset of the finite set \(m\), which is finite by item~3 of the finite-set
lemma; composing bijections, \(g[m]\) is finite. Every formula
\(\sigma\in\Delta\) is a string over the alphabet of \(\Lang^*\), a function
on some \(m\in\N\), so the set \(\sigma[m]\) of symbols occurring in it is
finite. Let \(b:k\to\Delta\) be a bijection, \(k\in\N\); the set
\(S=\bigcup_{i<k}\sigma_{b(i)}[\operatorname{dom}\sigma_{b(i)}]\) of
symbols occurring in members of \(\Delta\) is finite by item~4 of the
finite-set lemma. The map \(\alpha\mapsto c_\alpha\) restricted to
\(I_\Delta\) is an injection into \(S\); so \(I_\Delta\) is bijective with a
subset of \(S\), which is finite by item~3 of the finite-set lemma, and
\(I_\Delta\) is finite. \Cref{thm:p02-w02b-p02-029-finite-injections} gives an
injection \(\iota:I_\Delta\to D\) for every infinite \(D\).

This is the finite-satisfiability step of the proof of \cref{thm:upward-ls}:
given an infinite model \(\mathcal M_0\) of \(T\) with \(D=\dom\mathcal M_0\)
and a fixed \(d_0\in D\), interpret \(c_\alpha\) as \(\iota(\alpha)\) for
\(\alpha\in I_\Delta\) and as \(d_0\) otherwise. Every sentence
\(c_\alpha\ne c_\beta\) in \(\Delta\) has \(\alpha,\beta\in I_\Delta\) with
\(\alpha\neq\beta\), so it is satisfied because \(\iota\) is injective; the
members of \(\Delta\) from \(T\) are \(\Lang\)-sentences and are satisfied
because the expansion has \(\mathcal M_0\) as its \(\Lang\)-reduct.

\emph{Item~3.} For \(\alpha\neq\beta\) below \(\kappa\), item~3 of the ordinal
lemma gives \(\alpha<\beta\) or \(\beta<\alpha\), and the corresponding
sentence gives \(c_\alpha^{\mathcal M^*}\neq c_\beta^{\mathcal M^*}\). So
\(\alpha\mapsto c_\alpha^{\mathcal M^*}\) is injective; it maps onto \(A\)
by the definition of \(A\). As \(\kappa\) is a cardinal bijective with
\(A\), item~2 of cardinal assignment gives \(\abs A=\kappa\). By item~1,
\(\abs{\Lang^*}=\kappa\), and \(\aleph_0\leq\kappa\). The maximum of the
three is one of them, hence \(\leq\kappa\), and it is \(\geq\abs A=\kappa\);
as \(\leq\) on ordinals is inclusion, it equals \(\kappa\).

\emph{Item~4.} The bijection \(\kappa\to A\) followed by the inclusion
\(A\subseteq N\) is an injection \(\kappa\to N\), and \(\abs N\leq\kappa\) is
an injection \(N\to\kappa\). CSB gives a bijection, and item~2 of cardinal
assignment gives \(\abs N=\kappa\).
\end{proof}
\end{annexitem}

\begin{annexitem}{annex:p02-w02b-p02-029-terms-needed}{thm:p02-w02b-p02-029-terms-needed}{Each term of the downward bound is needed}
\begin{proof}
\emph{Item~1.} The inclusion \(A\subseteq\dom\mathcal N\) is an injection;
item~3 of cardinal assignment converts it into
\(\abs A\leq\abs{\dom\mathcal N}\).

\emph{Item~2.} For \(n\geq2\) let \(\iota_n\) be the sentence of
\cref{ch:fol-theories} asserting that \(x_0,\ldots,x_{n-1}\) can be chosen
pairwise distinct. By \cref{thm:p02-w02b-p02-029-finite-injections} there is
an injection \(e:n\to\dom\mathcal M\); an assignment sending \(x_i\) to
\(e(i)\) for \(i<n\) satisfies every conjunct \(x_i\ne x_j\) with \(i<j\),
so \(\mathcal M\entails\iota_n\). The definition of elementary substructure
with no displayed variables gives \(\mathcal N\entails\iota_n\): there are
pairwise distinct \(a_0,\ldots,a_{n-1}\in\dom\mathcal N\), that is, an
injection \(n\to\dom\mathcal N\). Suppose \(\dom\mathcal N\) were finite,
with a bijection \(b:\dom\mathcal N\to m\), \(m\in\N\). As
\(\dom\mathcal N\neq\varnothing\), \(m\geq1\) and \(m+1\geq2\); composing an
injection \(m+1\to\dom\mathcal N\) with \(b\) gives an injection
\(m+1\to m\), which item~1 of \cref{thm:p01-w02-p01-008-pigeonhole}
excludes. So \(\dom\mathcal N\) is infinite, its cardinal is bijective with
no \(n\in\N\), and item~1 of the infinite-ordinal lemma gives
\(\N\subseteq\abs{\dom\mathcal N}\), that is,
\(\aleph_0\leq\abs{\dom\mathcal N}\).

\emph{Item~3.} A substructure gives every constant the same value
(\cref{sec:fol-cardinality-normalforms-01}), so
\(C^{\mathcal M}\subseteq\dom\mathcal N\), and the argument of item~1
applies.

\emph{Item~4.} Apply \cref{thm:p02-w02b-p02-029-fresh} with
\(\Gamma=\mathsf{Var}\cup T_{\mathrm{log}}\) to obtain pairwise distinct
\(c_\xi\), \(\xi<\kappa\), outside \(\Gamma\), and let \(\Lang_\kappa\) have
exactly these as constant symbols and no other symbols. The map
\(\xi\mapsto c_\xi\) is a bijection from the cardinal \(\kappa\) onto the
symbol set, so \(\abs{\Lang_\kappa}=\kappa\) by item~2 of cardinal
assignment. Let \(\mathcal M_\kappa\) have domain \(\kappa\), nonempty
because \(\kappa\) is infinite, and \(c_\xi^{\mathcal M_\kappa}=\xi\). If
\(\mathcal N\subseteq\mathcal M_\kappa\), then
\(\xi=c_\xi^{\mathcal M_\kappa}=c_\xi^{\mathcal N}\in\dom\mathcal N\) for every
\(\xi<\kappa\), so \(\kappa\subseteq\dom\mathcal N\subseteq\kappa\). Every
elementary substructure is a substructure.
\end{proof}
\end{annexitem}
